\section{Conclusión}

La literatura revisada confirma una evolución clara en la detección computacional de señales de salud mental en redes sociales: el campo ha pasado de enfoques basados en rasgos léxicos y clasificadores convencionales a sistemas profundos, modelos basados en transformers y LLMs con una capacidad predictiva superior. Sin embargo, este avance no resuelve por sí mismo la exigencia central del dominio: producir salidas útiles, justificables y clínicamente comprensibles. Desde la dimensión de efectividad predictiva, el campo ha progresado de forma acelerada. Desde la dimensión de interpretabilidad clínica, en cambio, el progreso ha sido más desigual. Por ello, los trabajos más valiosos para este proyecto no son únicamente los que reportan la mejor exactitud, sino aquellos que articulan el rendimiento con estructuras interpretables, como niveles de severidad, categorías DSM-5, racionales expertos o extracción explícita de evidencia \citep{gaur2018dsm5, gaur2019severity, bao2025redsm5, singh2024evidence}.

En síntesis, el estado del arte sugiere que la dirección más prometedora no consiste en escoger entre desempeño e interpretabilidad, sino en diseñar tareas, datos y modelos que hagan compatibles ambas metas. Reddit se mantiene como una fuente especialmente valiosa por su riqueza discursiva, pero los retos de sesgo lingüístico, generalización, anotación experta y validación clínica siguen abiertos. Para las fases posteriores, esto implica buscar un equilibrio explícito entre desempeño cuantitativo, evidencia textual, modularidad y diseño reproducible, de modo que la herramienta resultante no solo clasifique mejor, sino que también permita justificar y auditar mejor sus decisiones.
